\noindent\textbf{Resultado.} La forma regular es
\[
-(xy')'=\lambda\frac1x y,
\qquad p=x,\quad r=0,\quad w=\frac1x.
\]
El producto adecuado es \(\langle f,g\rangle_w=\int_1^e fg\,dx/x\). Para el par dado,
\[
\int_1^e\frac{uv}{x}\,dx=0,\qquad
\int_1^e uv\,dx
=-\frac{4\pi^2(e+1)}{(1+\pi^2)(1+9\pi^2)}\ne0.
\]
Con \(z=\log x\), la medida ponderada se convierte en \(dz\); la no ponderada conserva el factor \(e^z\).

\noindent\textit{Desarrollo.} Dividir la ecuación por \(x>0\) da
\[
xy''+y'+\lambda\frac{y}{x}=0.
\]
Los dos primeros términos forman \((xy')'\), por lo que
\[
-(xy')'=\lambda\frac1x y.
\]
En \([1,e]\), \(p=x\) es continuamente derivable y positivo, \(r=0\) es continuo y \(w=1/x\) es continuo y positivo. Los contornos son Dirichlet homogéneos, de modo que el problema es regular.

Para interpretar el peso, ponemos \(z=\log x\) y \(Y(z)=y(e^z)\). Entonces \(xy'=Y_z\) y \(x^2y''+xy'=Y_{zz}\). El intervalo pasa a \([0,1]\) y la ecuación a
\[
-Y_{zz}=\lambda Y,\qquad Y(0)=Y(1)=0.
\]
Los perfiles dados corresponden a \(\sin(\pi z)\) y \(\sin(2\pi z)\), con autovalores \(\pi^2\) y \(4\pi^2\). Como \(dx/x=dz\), su producto ponderado es
\[
\int_1^e\frac{u(x)v(x)}x\,dx
=\int_0^1\sin(\pi z)\sin(2\pi z)\,dz=0.
\]
Esta es la ortogonalidad asociada al peso del problema.

Sin peso, \(dx=e^z\,dz\), y queda
\[
I=\int_0^1e^z\sin(\pi z)\sin(2\pi z)\,dz
=\frac12\int_0^1e^z[\cos(\pi z)-\cos(3\pi z)]\,dz.
\]
La primitiva
\[
\int e^z\cos(az)\,dz
=\frac{e^z[\cos(az)+a\sin(az)]}{1+a^2}
\]
se comprueba por derivación. Para \(a=\pi\) o \(a=3\pi\), los valores de borde proporcionan \(-(e+1)/(1+a^2)\). Por tanto,
\[
I=-\frac{e+1}{2}\left(\frac1{1+\pi^2}-\frac1{1+9\pi^2}\right)
=-\frac{4\pi^2(e+1)}{(1+\pi^2)(1+9\pi^2)}<0.
\]
No hay contradicción entre ambos resultados: se han usado medidas distintas. El cambio de variable conserva la ortogonalidad ordinaria en \(z\) al incluir el peso \(1/x\), no al omitirlo.
